% secV9_scope.tex -- Section 8. CONTAINER (27 Oct 2026): text to be written by the author.
% Word budget: 450 words (JPART count: running text, footnotes and citations as printed).
% Outline, numbers and sources: referee_EE_2026-10/ESQUELETO.md, section 8. The archived text
% (_arquivo_texto_IA_2026-10-27/) is not a source.

\section{Scope and limitations}\label{sec:scope}

The analysis establishes how alternative measurement conventions change a legal compliance classification. It also documents what enters a published waiting-list count and how reliably the supplied application files cover first requests. It does not estimate the causal effect of childcare expansion, registration practices, or the COVID-19 pandemic. Eleven candidate identification designs failed their prespecified diagnostic criteria, including instrument-strength, placebo, and covariate-balance checks; Table~\ref{tab:designs} in Appendix B records these attempts, including designs related to \citet{soares2021ecd} and \citet{schneider2023reforma}. No causal estimate is retained.

The compliance grid covers operational rules supported by the available aggregate series, not every interpretation of the statute. In particular, alternative treatments of pending documents and refusals are not separately evaluated. The narrower waiting list is observed only from September 2019 \citep{sme_demanda_escolar}; the School Census municipal-and-partnered split used here ends in 2021 \citep{inep_censo_escolar_c}. ``Never compliant'' therefore refers to a specification's evaluable years. We apply the legal thresholds without tolerance: the published December lists of 595 in 2021 and 681 in 2022 \citep{sme_demanda_escolar} remain positive, however small relative to enrolment. We have not calculated compliance under a permissible margin of error. The principal population construction assumes a linearly interpolated scaling factor between the two Demographic Censuses; the underlying law supplies no demographic series. The vintages of the SEADE projections \citep{seade_idade_escolar} have not been individually reconciled, and age-specific post-enumeration coverage estimates for S\~ao Paulo were unavailable. National Census coverage errors \citep{agenciabrasil2024ppe} cannot be applied to the city.

The administrative evidence has separate limitations. Extraction dates vary without an identified reason. In 2013--2018, reconstructed stocks are 4--13\% below the published counts \citep{sme_demanda_escolar}; the residual discrepancy is unexplained. The delivered request record ends with requests in 2018, preventing verification of the subsequent changes in list composition or the proposed flow measure in 2019--2022. Its coverage is assessed against another municipal record, not against all potentially eligible children. Estimates of the coverage-related component of the earlier cohort trend depend on the reference used for imputation, ranging from 29\% to 43\%. For 2.3--5.5\% of entrants, a request is recorded no more than seven days before enrolment, limiting a strictly prospective interpretation.

Finally, earlier registration is compatible with several mechanisms, including information, administrative procedures, and changes in childcare options; none is identified here. The data omit families who never apply, intended starting dates, income, informal care, and reasons for refusing offers. All compliance calculations are citywide; district-level disparities are not assessed. These limitations matter for the interpretation of observed counts and for the design of a better indicator, rather than invalidating the separate exercise based on published aggregate series.
